\documentclass{article}
\usepackage[utf8]{inputenc}
\usepackage[T1]{fontenc}
\usepackage{geometry}
\geometry{
    a4paper,
    left=2.5cm,
    right=2.5cm,
    top=2.5cm,
    bottom=2.5cm
}
\usepackage{graphicx}
\usepackage{float}
\usepackage{hyperref}
\usepackage[english]{babel}
\usepackage{amsmath}
\usepackage{amssymb}
\usepackage{array}
\usepackage{booktabs}
\usepackage{caption}
\usepackage{subcaption}
\usepackage{enumitem}
\usepackage{fancyhdr}
\usepackage{titlesec}
\usepackage{url}
\usepackage{textcomp}
\usepackage{xcolor}
\usepackage{listings}
\usepackage{multirow}
\usepackage{longtable}
\usepackage{tabularx}
\usepackage{times}

% Set up page style
\pagestyle{fancy}
\fancyhf{}
\rfoot{\thepage}
\renewcommand{\headrulewidth}{0pt}

% Title formatting
\titleformat{\chapter}[display]{\normalfont\bfseries}{\chaptertitle\thechapter}{0pt}{\Huge}
\titlespacing*{\chapter}{0pt}{50pt}{40pt}

% Document info
\title{Template Document}
\author{Academic Support Agent}
\date{\today}

\begin{document}
\maketitle

\begin{frame}{}
    \centering
    \hspace*{-0.5cm}
    $\vcenter{\hbox{\includegraphics[width=1cm]{img/Ufu_logo.png}}}$
    $\vcenter{\resizebox{0.95\textwidth}{!}{
        \begin{tabular}{c}
             FEELT31109 - Enriquecimento Instrumental $\cdot$ 2019-1 $\cdot$ Relat rio Lab. n  XX  \\
             \hline 
        \end{tabular}
    }}$
\end{frame}

\section*{Equipe}

\begin{itemize}
    \item Brenda Almeida Cavalcanti 11111ECP001
    \item Douglas Cardoso Araujo 11111ECP095
\end{itemize}

\section{Introdu o}

Breve descri o do que era esperado para cada projeto.

\section{Descri o dos experimentos}

Relato de cada uma das montagens e das dificuldades encontradas.

Se precisar de imagens, eis um exemplo na Figura~\ref{fig:exemplo}.

\begin{figure}[h!]
\centering
\includegraphics[height=5cm]{img/Lab01a_SOS_Luminoso_bb.png}
\caption{Exemplo de imagem}
\label{fig:exemplo}
\end{figure}

Se precisar de tabelas, eis um exemplo na Tabela~\ref{tab:exemplo}. Pelo site \url{https://tablesgenerator.com} voc  consegue montar tabelas do \LaTeX{}.

\begin{table}[htb]
\caption{Exemplo de tabela}
\centering
\label{tab:exemplo}
    \begin{tabular}{rcc}
        \hline \hline
        $\alpha$ & $\Gamma_{meas}$ & $\Gamma_{theor}$ \\
        \hline
        0 & 3.4 & 3.2 \\
        2 & 5.6 & 5.3 \\
        4 & 10.5 & 10.8 \\
        6 & 24.6 & 22.5 \\
        \hline
    \end{tabular}
\end{table}

\section{Resultados}

Relato dos resultados obtidos em cada um dos experimentos.

\paragraph{Opcional:} Descri o das modifica es feitas ao(s) projeto(s) e resultados
obtidos a partir das mesmas.
\end{document}